\documentclass[../../main.tex]{subfiles}
\begin{document}

\begin{center}
  {\fontsize{48pt}{58pt}\selectfont \textbf{1970}\normalsize\textbf{年度}}
\end{center}

\begin{table}[htb]
  \centering
  \caption{出題テーマと難易度}
  \begin{tabular}{|c|p{9cm}|c|}
    \hline
    \textbf{問題番号} & \textbf{テーマ} & \textbf{難易度} \\
    \hline
    第1問 & 複素数の累乗と積の値 & \\
    \hline
    第2問 & 硬貨投げによる一次元の移動 & \\
    \hline
    第3問 & 二人の移動と出会うための速さ & \\
    \hline
    第4問 & 水槽の穴から出る水の軌道 & \\
    \hline
  \end{tabular}
\end{table}

\begin{table}[htb]
  \centering
  \caption{解答}
  \renewcommand{\arraystretch}{1.4}
  \begin{tabular}{|c|c|p{9cm}|}
    \hline
    \textbf{問題番号} & \textbf{小問} & \textbf{答え} \\
    \hline
    第1問 & (1) & $6$個 \\
    \cline{2-3}
    & (2) & $n\equiv\pm1\pmod{6}$なら$1$，それ以外なら$0$ \\
    \hline
    第2問 & (1) & $x=0$で$\dfrac{3}{8}$，$x=\pm2$で$\dfrac{1}{4}$，$x=\pm4$で$\dfrac{1}{16}$，その他は$0$ \\
    \cline{2-3}
    & (2) & $x=n-2$で$\dfrac{n}{2^n}$，$x=n-4$で$\dfrac{n(n-1)}{2^{n+1}}$ \\
    \hline
    第3問 & -- & $u$の最小値は$\dfrac{15}{4}\,\mathrm{m/s}$ \\
    \hline
    第4問 & (1) & 穴の真下から$2\sqrt{h(a-h)}$だけ離れた点 \\
    \cline{2-3}
    & (2) & 水面からの深さを$h=\dfrac{a}{2}$とする \\
    \cline{2-3}
    & (3) & 通過範囲は$0\le X\le a$，$0\le Y\le a-X$（図示は解答参照） \\
    \hline
  \end{tabular}
\end{table}

\paragraph{試験形態}
4問．試験時間は要確認．

\end{document}
